% Generador de dorsos: modelo geométrico y de simetría.
%
% Compilación (desde docs/methodology/):
%   pdflatex model.tex
%   pdflatex model.tex      % segunda pasada para las referencias internas
% Paquetes: amsmath, amssymb, amsthm, booktabs, geometry, hyperref y babel. El
% español se carga con \babelprovide, igual que en docs/methodology/brisca.tex
% del repositorio.

\documentclass[11pt,a4paper]{article}
\usepackage[utf8]{inputenc}
\usepackage[T1]{fontenc}
\usepackage{babel}
\babelprovide[import, main]{spanish}
\usepackage{amsmath,amssymb,amsthm}
\usepackage{booktabs}
\usepackage[margin=2.5cm]{geometry}
\usepackage{hyperref}

\theoremstyle{definition}
\newtheorem{definicion}{Definición}[section]
\newtheorem{supuesto}{Supuesto}[section]
\theoremstyle{plain}
\newtheorem{proposicion}{Proposición}[section]
\newtheorem{lema}{Lema}[section]
\theoremstyle{remark}
\newtheorem{observacion}{Observación}[section]

\newcommand{\R}{\mathbb{R}}
\newcommand{\Z}{\mathbb{Z}}
\newcommand{\Klein}{\mathcal{K}}

\title{Generador de dorsos: modelo geométrico y de simetría}
\author{Juegos reunidos}
\date{}

\begin{document}
\maketitle

\section{Convenciones y paso de milímetros a píxeles}

\begin{definicion}[Carta]
Una carta es el rectángulo $R=[0,W]\times[0,H]\subset\R^2$, en milímetros, con el
origen arriba a la izquierda y el eje $y$ hacia abajo. Su centro es
$c=(W/2,\,H/2)$. Los formatos son $(W,H)=(63,88)$ (póker) y $(61{,}5,\,95)$
(baraja española).
\end{definicion}

Sea $d$ la resolución en puntos por pulgada y $S$ el factor de sobremuestreo
(par). El tamaño final en píxeles es
\[
  p_x=\operatorname{round}\!\Bigl(\tfrac{W}{25{,}4}\,d\Bigr),\qquad
  p_y=\operatorname{round}\!\Bigl(\tfrac{H}{25{,}4}\,d\Bigr).
\]
Se dibuja en un lienzo de $Sp_x\times Sp_y$ píxeles con escalas por eje
\[
  k_x=\frac{Sp_x}{W},\qquad k_y=\frac{Sp_y}{H},\qquad (x,y)\mapsto(k_xx,\;k_yy).
\]

\begin{proposicion}\label{prop:centro}
El centro de la carta se transforma en el centro del lienzo, $(Sp_x/2,\,Sp_y/2)$,
y el error relativo de escala respecto de $S\,d/25{,}4$ es a lo sumo $1/(2p_x)$
en $x$ (análogo en $y$).
\end{proposicion}
\begin{proof}
$k_x\,W/2=Sp_x/2$ por definición. Además $p_x=\tfrac{W d}{25{,}4}+\varepsilon$
con $|\varepsilon|\le\tfrac12$, luego $k_x=\tfrac{S d}{25{,}4}\bigl(1+\tfrac{25{,}4\,\varepsilon}{Wd}\bigr)$
y $\bigl|\tfrac{25{,}4\,\varepsilon}{Wd}\bigr|\le\tfrac{1}{2p_x}\,(1+o(1))$.
\end{proof}

La PNG guarda $d$ en su cabecera, de modo que $p_x/d\cdot25{,}4$ reproduce $W$ con
error a lo sumo $12{,}7/d$~mm (0,04~mm a 300~ppp). El SVG usa directamente
milímetros: \texttt{width="$W$mm"} y \texttt{viewBox="0 0 $W$ $H$"}.

El grosor de una línea de $w$~mm se convierte en $\operatorname{round}(w\,(k_x+k_y)/2)$
píxeles del lienzo sobremuestreado, con un mínimo de 1.

\section{Simetría como acción de un grupo}

\subsection{El grupo}
Sean $m_x(x,y)=(W-x,\,y)$ (espejo respecto del eje vertical central),
$m_y(x,y)=(x,\,H-y)$ (espejo respecto del horizontal) y
$r(x,y)=(W-x,\,H-y)$ (giro de $180^\circ$ alrededor de $c$).

\begin{proposicion}
$\Klein=\{e,m_x,m_y,r\}$ es un grupo de isometrías de $R$ isomorfo al de Klein
$\Z_2\times\Z_2$, con $m_x\circ m_y=r$.
\end{proposicion}
\begin{proof}
Cada elemento es una involución: $m_x^2(x,y)=(W-(W-x),y)=(x,y)$, y análogamente
$m_y^2=r^2=e$. Además $m_x(m_y(x,y))=m_x(x,H-y)=(W-x,H-y)=r(x,y)$, y $m_y\circ m_x=r$
por el mismo cálculo, luego el producto es conmutativo y cerrado. Todo elemento
distinto de $e$ tiene orden 2 y hay cuatro, que es la definición de $\Z_2\times\Z_2$.
\end{proof}

Los tres subgrupos que ofrece el generador son $\{e,m_x\}$ (simetría
\emph{bilateral}), $\{e,r\}$ (giro de $180^\circ$) y $\Klein$ (ambas; implica la
simetría respecto del eje horizontal, porque $m_y=r\circ m_x$).

\subsection{Dominio fundamental y construcción}
\begin{definicion}
Para un subgrupo $G\le\Klein$, un \emph{dominio fundamental} $D\subset R$ cumple
$R=\bigcup_{g\in G}g(D)$ con intersecciones de medida nula.
\end{definicion}

Se toma $D=[0,W/2]\times[0,H]$ para $\{e,m_x\}$; $D=[0,W]\times[0,H/2]$ para
$\{e,r\}$; y $D=[0,W/2]\times[0,H/2]$ para $\Klein$. En los tres casos las
imágenes $g(D)$ son rectángulos que teselan $R$.

Cada módulo dibuja la carta completa, lo que da una función
$f:R\to\text{colores}$ (la imagen sin simetrizar). La imagen final es
\[
  F(p)=f\bigl(g^{-1}p\bigr)\quad\text{si }p\in g(D).
\]

\begin{proposicion}
$F$ es $G$-invariante: $F(hp)=F(p)$ para todo $h\in G$ y $p\in R$ (salvo en los
bordes de medida nula, donde $F$ toma el valor de cualquiera de los dominios que
tocan).
\end{proposicion}
\begin{proof}
Si $p\in g(D)$, entonces $hp\in(hg)(D)$ y, por definición,
$F(hp)=f\bigl((hg)^{-1}hp\bigr)=f(g^{-1}p)=F(p)$.
\end{proof}

\begin{observacion}
La construcción no exige que $f$ sea simétrica: solo se usa la restricción
$f|_D$. Por eso los marcos, medallones y esquinas pueden dibujarse sin conocer la
simetría, y una semilla aleatoria por celda produce igualmente un diseño
$G$-invariante.
\end{observacion}

\subsection{Realización exacta sobre píxeles}
Si el lienzo mide $2a\times2b$ píxeles (par en ambas dimensiones), los giros y
espejos del generador son las transposiciones de imagen exactas: el píxel
$(u,v)$ va a $(2a-1-u,\,v)$ para $m_x$, a $(u,\,2b-1-v)$ para $m_y$ y a
$(2a-1-u,\,2b-1-v)$ para $r$. Como $S$ es par, $Sp_x$ y $Sp_y$ lo son, y por la
proposición~\ref{prop:centro} el eje de simetría del dibujo coincide con el del
lienzo. El reescalado por $S$ es un filtro de caja sobre bloques $S\times S$
alineados; los bloques se transforman en bloques, así que la simetría se conserva
píxel a píxel tras reducir. Las pruebas automáticas lo comprueban con diferencia
cero entre la imagen y sus transpuestas.

En SVG se dibuja la carta una vez en un \texttt{<g>} y se añaden, por cada
$g\ne e$, un \texttt{<use>} con la matriz de $g$ recortado al rectángulo
$g(D)$. La matriz de $g(x,y)=(ax+by+e,\,cx+dy+f)$ se escribe
\texttt{matrix(a c b d e f)}, porque SVG ordena las columnas.

\section{Retículas}

\begin{definicion}[Retícula centrada]
Sea el campo $[x_0,x_1]\times[y_0,y_1]$ de ancho $w$ y alto $h$, con centro
$(c_x,c_y)$, y $s_0=\text{cell\_mm}\times\text{scale}$ el tamaño nominal de celda.
\begin{align*}
  n_x&=\operatorname{impar}\!\bigl(\operatorname{round}(w/s_0)\bigr), & s&=w/n_x,\\
  n_y&=\operatorname{impar}\!\bigl(\operatorname{round}(h/s)\bigr),
\end{align*}
donde $\operatorname{impar}(n)$ es el impar más cercano a $w/s_0$ (o $h/s$) entre
$n$ y $n\pm1$. Los nodos son $v_{ij}=(c_x+(i-\tfrac{n_x}{2})s,\;c_y+(j-\tfrac{n_y}{2})s)$
y el centro de la celda $(i,j)$ es $v_{i+\frac12,\,j+\frac12}$.
\end{definicion}

Así, un número entero de celdas cubre exactamente el ancho, y en vertical sobran
$(h-n_ys)/2$ arriba y abajo: las celdas de los extremos quedan cortadas por igual,
como en las barajas reales. El recorte al interior del marco lo hace un polígono.

\begin{proposicion}\label{prop:paridad}
Si $n_x$ y $n_y$ son impares, el espejo $m_x$ lleva la celda $(i,j)$ a la
$(n_x-1-i,\,j)$ y el giro $r$ a la $(n_x-1-i,\,n_y-1-j)$, y ambas conservan la
paridad de $i+j$.
\end{proposicion}
\begin{proof}
El nodo $v_{ij}$ tiene abscisa $c_x+(i-n_x/2)s$; su imagen por $m_x$ es
$2c_x-\bigl(c_x+(i-n_x/2)s\bigr)=c_x+(n_x-i-n_x/2)s$, es decir el nodo $n_x-i$. El
centro de celda $i+\tfrac12$ pasa a $n_x-i-\tfrac12=(n_x-1-i)+\tfrac12$. La paridad
de $n_x-1-i$ es la de $i$ si $n_x-1$ es par, esto es, si $n_x$ es impar. Lo mismo
en $j$; el giro compone ambos.
\end{proof}

Por eso las variantes que alternan color o motivo por paridad (\emph{damero},
\emph{mosaico}) usan celdas impares: con un número par la alternancia se
invertiría en el eje y aparecerían dos celdas iguales contiguas justo en la
costura.

\subsection{Malla de rombos}
Para los rombos se usa el paso $s/2$ y solo los nodos con $i+j$ par, con
$n_x,n_y$ pares. Un rombo de semidiagonal $s/2$ centrado en $v_{ij}$ es
$\{|x-x_{ij}|+|y-y_{ij}|\le s/2\}$. Dos rombos vecinos $(i,j)$ y $(i+1,j\pm1)$
distan $(s/2,\pm s/2)$ y sus bordes comparten exactamente el punto medio
$v_{ij}+(s/4,\pm s/4)$, que cumple $|s/4|+|s/4|=s/2$; el área de un rombo es
$2(s/2)^2=s^2/2$ y la densidad de nodos con $i+j$ par es $1/(2(s/2)^2)=2/s^2$, luego
la malla cubre el plano sin huecos ni solapes ($s^2/2\cdot2/s^2=1$).

Dos rombos comparten arista si y solo si $(\Delta i,\Delta j)=(\pm1,\pm1)$. La
función $c(i,j)=i\bmod2$ cambia en todos esos desplazamientos, luego es un
coloreado propio con dos colores (el damero de la malla girada $45^\circ$). El
coloreado $\lfloor(i-j)/2\rfloor\bmod2$ \emph{no} lo es: da rayas diagonales, y
por eso el generador usa $i\bmod2$.

\section{Curvas}

\subsection{Unión de círculos}
Los medallones de roseta y cuatrilóbulo son la unión de círculos $C_k$ de centro
$q_k$ y radio $\rho_k$ con $|q_k|\le\rho_k$ (contienen o rozan el origen).

\begin{proposicion}
La unión es estrellada respecto del origen y su función radial es
\[
  \varrho(\theta)=\max_k\Bigl[\,u\cdot q_k+\sqrt{(u\cdot q_k)^2-|q_k|^2+\rho_k^2}\,\Bigr],
  \qquad u=(\cos\theta,\sin\theta).
\]
\end{proposicion}
\begin{proof}
El rayo $tu$, $t\ge0$, corta a $C_k$ donde $|tu-q_k|^2=\rho_k^2$, es decir
$t^2-2t\,(u\cdot q_k)+|q_k|^2-\rho_k^2=0$. El producto de raíces es
$|q_k|^2-\rho_k^2\le0$ y el discriminante reducido es $(u\cdot q_k)^2+\rho_k^2-|q_k|^2\ge0$,
así que hay una raíz no negativa, $t_k^+=u\cdot q_k+\sqrt{\cdot}$. Cada disco es
estrellado respecto del origen porque es convexo y lo contiene; una unión de
conjuntos estrellados respecto de un punto lo es, con radio el máximo.
\end{proof}

\paragraph{Lóbulos de la roseta.} Con $n$ lóbulos de radio $\rho$ a distancia
$d=1-\rho$ del centro, dos lóbulos contiguos (centros a distancia $2d\sin(\pi/n)$)
son tangentes cuando $2d\sin(\pi/n)=2\rho$, esto es
$\rho=\dfrac{\sin(\pi/n)}{1+\sin(\pi/n)}$. El generador usa $1{,}25\rho$ para que se
solapen; el contorno alcanza entonces el radio $1+0{,}25\rho$, algo mayor que el
nominal $1$.

\paragraph{Cuatrilóbulo.} Círculos de radio $1/2$ en $(\pm\tfrac12,0)$ y
$(0,\pm\tfrac12)$: extensión 1 sobre los ejes; el origen es un punto de la
frontera de cada uno ($|q|=\rho$) pero la unión lo rodea, y la fórmula sigue
siendo válida.

\subsection{Rosetón de guilloché}
Con $k$ lóbulos, profundidad $\delta\in(0,1)$ y $m$ curvas,
\[
  r_c(\theta)=R\Bigl[1-\tfrac{\delta}{2}\bigl(1+\cos k(\theta-\varphi_c)\bigr)\Bigr],
  \qquad \varphi_c=\frac{2\pi c}{mk},\quad c=0,\dots,m-1.
\]
Cada curva oscila entre $R(1-\delta)$ y $R$ con $k$ lóbulos, y las $m$ copias
están giradas una fracción del período $2\pi/k$. Dos curvas $c\ne c'$ se cortan
donde $\cos k(\theta-\varphi_c)=\cos k(\theta-\varphi_{c'})$; como
$\cos A=\cos B\iff A=\pm B+2\pi j$ y el caso $A=B$ exige $\varphi_c=\varphi_{c'}$,
queda $A=-B$, o sea $\theta=\tfrac{\varphi_c+\varphi_{c'}}2+\tfrac{\pi j}{k}$,
$j=0,\dots,2k-1$: $2k$ cruces por pareja, que forman el moaré característico. Se
dibujan dos anillos ($R=0{,}47s$, $\delta=0{,}45$ y $R=0{,}26s$, $\delta=0{,}5$).

\paragraph{Haces sinusoidales.} Cada fila usa $y=y_r+A\cos\!\bigl(\tfrac{2\pi(x-c_x)}{\lambda}+\pi\tfrac{c}{m}\bigr)$
con $c=0,\dots,m-1$. La fase reparte las curvas en $[0,\pi)$, de modo que se
cruzan en $x$ donde $\cos(\alpha+\pi c/m)=\cos(\alpha+\pi c'/m)$.

\subsection{Voluta}
La espiral de Arquímedes de centro $q$ es $r(t)=r_0+(r_1-r_0)t$,
$\theta(t)=\theta_0+2\pi Nt$, $t\in[0,1]$, con $N$ vueltas. La separación
radial entre vueltas consecutivas es $\Delta r=|r_1-r_0|/N$; para que el trazo de
ancho $w$ no se funda consigo mismo hace falta $\Delta r>w$. Con $s$ la celda,
$r_0=0{,}19s$, $r_1=0{,}035s$, $N=1{,}5$ da $\Delta r\approx0{,}10s$.
La voluta del primer cuadrante se refleja en los otros tres con $m_x$, $m_y$
y $r$ respecto del centro de la celda, así que la celda tiene simetría $\Klein$.

\subsection{Hoja (lente)}
Entre $p$ y $q$, con normal unitaria $\nu$ y flecha $b$, los dos arcos se aproximan por
la parábola $h(t)=4b\,t(1-t)$, $t\in[0,1]$: $\gamma_\pm(t)=p+t(q-p)\pm h(t)\nu$.
La flecha máxima es $h(\tfrac12)=b$ y la curva es tangente en $p$ y $q$ a
$\pm4b\nu$.

\subsection{Estrella cóncava (ogivas)}
En una celda de lado $s$ se centran cuatro círculos de radio $s/2$ en sus
esquinas. Los centros contiguos distan $s=2\cdot\tfrac s2$, luego son tangentes en
los puntos medios de los lados. La región de la celda no cubierta por los cuatro
cuartos de disco es una estrella de cuatro puntas con área
$s^2-4\cdot\tfrac14\pi(s/2)^2=s^2(1-\pi/4)\approx0{,}215\,s^2$.

\subsection{Escamas y bandas}
Las escamas son círculos de radio $s/2$ centrados en los nodos de la malla de
rombos; se pintan por filas de arriba abajo (algoritmo del pintor) y cada fila
oculta la parte baja de la anterior. Las bandas sinuosas son
$y=y_r+A\cos\bigl(2\pi(x-c_x)/\lambda\bigr)$ con $\lambda=2s$ y $A=0{,}22s$; al ser
pares en $x-c_x$ son simétricas respecto del eje vertical sin más.

\section{Marcos con banda}
Sea $O$ el rectángulo exterior del marco, $t$ el ancho de banda e
$I=O$ contraído $t$ por cada lado. Las cuatro bandas se recorren en coordenadas
locales $(u,v)$, con $u\in[0,L]$ a lo largo del lado y $v\in[0,t]$ desde el borde
exterior. Las transformaciones a página son rotaciones múltiplos de
$\pi/2$:
\begin{align*}
  A_\text{arriba}(u,v)&=(x_0+t+u,\;y_0+v), & L&=w-2t,\\
  A_\text{derecha}(u,v)&=(x_1-v,\;y_0+t+u), & L&=h-2t,\\
  A_\text{abajo}(u,v)&=(x_1-t-u,\;y_1-v), & L&=w-2t,\\
  A_\text{izquierda}(u,v)&=(x_0+v,\;y_1-t-u), & L&=h-2t,
\end{align*}
con $O=[x_0,x_1]\times[y_0,y_1]$. Cada una tiene matriz ortogonal de determinante
$+1$ (p.\,ej.\ la de la derecha: $\left(\begin{smallmatrix}0&-1\\1&0\end{smallmatrix}\right)$),
luego es un giro. El lado se divide en $n=\max\bigl(1,\operatorname{round}(L/(t\,\mu))\bigr)$
unidades de longitud $L/n$, con $\mu$ la razón de aspecto de la unidad. Si el motivo
es simétrico respecto del eje $u=L/(2n)$ de su unidad, el lado lo es respecto de su
punto medio. En las esquinas hay un bloque $t\times t$ independiente.

\section{Ornamentos de esquina}
Se diseñan para la esquina superior izquierda $(x_0,y_0)$ del interior del marco,
con $u$ hacia dentro y $v$ hacia abajo, y se colocan con
\begin{align*}
  M_\text{sup-izq}(u,v)&=(x_0+u,\;y_0+v), & M_\text{sup-der}(u,v)&=(x_1-u,\;y_0+v),\\
  M_\text{inf-der}(u,v)&=(x_1-u,\;y_1-v), & M_\text{inf-izq}(u,v)&=(x_0+u,\;y_1-v),
\end{align*}
es decir, con las reflexiones de $\Klein$ alrededor de los ejes del interior.

\section{Reproducibilidad y randomización}
\begin{supuesto}
La semilla de Python \texttt{random.Random(str)} es determinista entre plataformas
(se deriva de la cadena por SHA-512).
\end{supuesto}
Cada módulo y cada celda obtienen un generador propio con la cadena
``\texttt{seed/\allowbreak ranura/\allowbreak i/\allowbreak j}''. Los flujos son independientes: activar o desactivar
un módulo no altera las decisiones de los demás, y las decisiones de una celda no
dependen del orden de iteración. La randomización $\operatorname{randomize}(b,\sigma,\Lambda)$
construye un diseño a partir de $\sigma$ y sobrescribe cada ruta $\lambda\in\Lambda$
con el valor que tiene en la configuración base $b$; por construcción todo lo
bloqueado se conserva y lo demás es función de $\sigma$.

\paragraph{Limitaciones.} (i) La simetría es exacta a nivel de píxel en el PNG; en el
SVG depende del antialiasing del visor en la costura central. (ii) El diseño es
determinista dada la configuración y la versión del generador: cambiar
las constantes de un módulo cambia la imagen aunque la semilla sea la misma.
(iii) No se ha medido la calidad visual con ningún criterio objetivo: los
intervalos de randomización se eligieron a ojo sobre hojas de muestras.

\end{document}
